\section{AMC 10 10A 2005}
        \begin{enumerate}
        	\item (\textit{AMC 10 10A 2005}) While eating out, Mike and Joe each tipped their server $\$2$. Mike tipped $10\%$ of his bill and Joe tipped $20\%$ of his bill. What was the difference, in dollars, between their bills? 


 $\textbf{(A) } 2\qquad \textbf{(B) } 4\qquad \textbf{(C) } 5\qquad \textbf{(D) } 10\qquad \textbf{(E) } 20$


	\item (\textit{AMC 10 10A 2005}) For each pair of real numbers $a \neq b$, define the operation $\star$ as


 
 \[(a \star b) = \frac{a+b}{a-b}.\]


 What is the value of $\left(\left(1 \star 2\right) \star 3\right)$?


 $\textbf{(A) } -\frac{2}{3}\qquad \textbf{(B) } -\frac{1}{5}\qquad \textbf{(C) } 0\qquad \textbf{(D) } \frac{1}{2}\qquad \textbf{(E) } \text{This value is not defined.}$


	\item (\textit{AMC 10 10A 2005}) The equations $2x + 7 = 3$ and $bx - 10 = -2$ have the same solution $x$. What is the value of $b$? 


 $\textbf{(A) } -8\qquad \textbf{(B) } -4\qquad \textbf{(C) } -2\qquad \textbf{(D) } 4\qquad \textbf{(E) } 8$


	\item (\textit{AMC 10 10A 2005}) A rectangle with a diagonal of length $x$ is twice as long as it is wide. What is the area of the rectangle? 


 $\textbf{(A) } \frac{1}{4}x^2\qquad \textbf{(B) } \frac{2}{5}x^2\qquad \textbf{(C) } \frac{1}{2}x^2\qquad \textbf{(D) } x^2\qquad \textbf{(E) } \frac{3}{2}x^2$


	\item (\textit{AMC 10 10A 2005}) A store normally sells windows at $\$100$ each. This week the store is offering one free window for each purchase of four. Dave needs seven windows and Doug needs eight windows. How many dollars will they save if they purchase the windows together rather than separately?


 $\textbf{(A) } 100\qquad \textbf{(B) } 200\qquad \textbf{(C) } 300\qquad \textbf{(D) } 400\qquad \textbf{(E) } 500$


	\item (\textit{AMC 10 10A 2005}) The average (mean) of $20$ numbers is $30$, and the average of $30$ other numbers is $20$. What is the average of all $50$ numbers?


 $\textbf{(A) } 23\qquad \textbf{(B) } 24\qquad \textbf{(C) } 25\qquad \textbf{(D) } 26\qquad \textbf{(E) } 27$


	\item (\textit{AMC 10 10A 2005}) Josh and Mike live $13$ miles apart. Josh started to ride his bicycle toward Mike's house. A little later Mike started to ride his bicycle toward Josh's house. When they met, Josh had ridden for twice the length of time as Mike and at four-fifths of Mike's rate. How many miles had Mike ridden when they met? 


 $\textbf{(A) } 4\qquad \textbf{(B) } 5\qquad \textbf{(C) } 6\qquad \textbf{(D) } 7\qquad \textbf{(E) } 8$


	\item (\textit{AMC 10 10A 2005}) In the figure, the length of side $AB$ of square $ABCD$ is $\sqrt{50}$, $E$ is between $B$ and $H$, and $BE = 1$. What is the area of the inner square $EFGH$?


 
                    \begin{figure}[H]
                        \centering
                        \includegraphics[width=0.5\textwidth]{Problem/AMC_10_10A_2005_8_Problem_Figure_1.png}
                    \end{figure}
                    
$\textbf{(A) } 25\qquad\textbf{(B) } 32\qquad\textbf{(C) } 36\qquad\textbf{(D) } 40\qquad\textbf{(E) } 42$


	\item (\textit{AMC 10 10A 2005}) Three tiles are marked $X$ and two other tiles are marked $O$. The five tiles are randomly arranged in a row. What is the probability that the arrangement reads $XOXOX$?


 $\textbf{(A) } \frac{1}{12}\qquad \textbf{(B) } \frac{1}{10}\qquad \textbf{(C) } \frac{1}{6}\qquad \textbf{(D) } \frac{1}{4}\qquad \textbf{(E) } \frac{1}{3}$


	\item (\textit{AMC 10 10A 2005}) There are two values of $a$ for which the equation $4x^2 + ax + 8x + 9 = 0$ has only one solution for $x$. What is the sum of those values of $a$?


 $\textbf{(A) } -16\qquad \textbf{(B) } -8\qquad \textbf{(C) } 0\qquad \textbf{(D) } 8\qquad \textbf{(E) } 20$


	\item (\textit{AMC 10 10A 2005}) A wooden cube $n$ units on a side is painted red on all six faces and then cut into $n^3$ unit cubes. Exactly one-fourth of the total number of faces of the unit cubes are red. What is $n$?


 $\textbf{(A) } 3\qquad \textbf{(B) } 4\qquad \textbf{(C) } 5\qquad \textbf{(D) } 6\qquad \textbf{(E) } 7$


	\item (\textit{AMC 10 10A 2005}) The figure shown is called a \textit{trefoil} and is constructed by drawing circular sectors about sides of the congruent equilateral triangles. What is the area of a trefoil whose horizontal base has length $2$?


 
                    \begin{figure}[H]
                        \centering
                        \includegraphics[width=0.5\textwidth]{Problem/AMC_10_10A_2005_12_Problem_Figure_1.png}
                    \end{figure}
                    
$\textbf{(A) } \frac{1}{3}\pi+\frac{\sqrt{3}}{2}\qquad \textbf{(B) } \frac{2}{3}\pi\qquad \textbf{(C) } \frac{2}{3}\pi+\frac{\sqrt{3}}{4}\qquad \textbf{(D) } \frac{2}{3}\pi+\frac{\sqrt{3}}{3}\qquad \textbf{(E) } \frac{2}{3}\pi+\frac{\sqrt{3}}{2}$


	\item (\textit{AMC 10 10A 2005}) How many positive integers $n$ satisfy the following condition:


 
 \[\left(130n\right)^{50} > n^{100} > 2^{200} \ \text{?}\]
$\textbf{(A) } 0\qquad \textbf{(B) } 7\qquad \textbf{(C) } 12\qquad \textbf{(D) } 65\qquad \textbf{(E) } 125$


	\item (\textit{AMC 10 10A 2005}) How many three-digit numbers satisfy the property that the middle digit is the average of the first and the last digits? 


 $\textbf{(A) } 41\qquad \textbf{(B) } 42\qquad \textbf{(C) } 43\qquad \textbf{(D) } 44\qquad \textbf{(E) } 45$


	\item (\textit{AMC 10 10A 2005}) How many positive cubes divide $3! \cdot 5! \cdot 7!$?


 $\textbf{(A) } 2\qquad \textbf{(B) } 3\qquad \textbf{(C) } 4\qquad \textbf{(D) } 5\qquad \textbf{(E) } 6$


	\item (\textit{AMC 10 10A 2005}) The sum of the digits of a two-digit number is subtracted from the number. The units digit of the result is $6$. How many two-digit numbers have this property? 


 $\textbf{(A) } 5\qquad \textbf{(B) } 7\qquad \textbf{(C) } 9\qquad \textbf{(D) } 10\qquad \textbf{(E) } 19$


	\item (\textit{AMC 10 10A 2005}) In the five-sided star shown, the letters $A$, $B$, $C$, $D$, and $E$ are replaced by the numbers $3$, $5$, $6$, $7$, and $9$, although not necessarily in this order. The sums of the numbers at the ends of the line segments $\overline{AB}$, $\overline{BC}$, $\overline{CD}$, $\overline{DE}$, and $\overline{EA}$ form an arithmetic sequence, although not necessarily in this order. What is the middle term of the arithmetic sequence? 


 
                    \begin{figure}[H]
                        \centering
                        \includegraphics[width=0.5\textwidth]{Problem/AMC_10_10A_2005_17_Problem_Figure_1.png}
                    \end{figure}
                    
$\textbf{(A) } 9\qquad \textbf{(B) } 10\qquad \textbf{(C) } 11\qquad \textbf{(D) } 12\qquad \textbf{(E) } 13$


	\item (\textit{AMC 10 10A 2005}) Team $A$ and team $B$ play a series. The first team to win three games wins the series. Each team is equally likely to win each game, there are no ties, and the outcomes of the individual games are independent. If team $B$ wins the second game and team $A$ wins the series, what is the probability that team $B$ wins the first game? 


 $\textbf{(A) } \frac{1}{5}\qquad \textbf{(B) } \frac{1}{4}\qquad \textbf{(C) } \frac{1}{3}\qquad \textbf{(D) } \frac{1}{2}\qquad \textbf{(E) } \frac{2}{3}$


	\item (\textit{AMC 10 10A 2005}) Three one-inch squares are placed with their bases on a line. The center square is lifted out and rotated $45^{\circ}$, as shown. Then it is centered and lowered into its original location until it touches both of the adjoining squares. How many inches is the point $B$ from the line on which the bases of the original squares were placed?


 
                    \begin{figure}[H]
                        \centering
                        \includegraphics[width=0.5\textwidth]{Problem/AMC_10_10A_2005_19_Problem_Figure_1.png}
                    \end{figure}
                    
$\textbf{(A) } 1\qquad\textbf{(B) } \sqrt{2}\qquad\textbf{(C) } \frac{3}{2}\qquad\textbf{(D) } \sqrt{2}+\frac{1}{2}\qquad\textbf{(E) } 2$


	\item (\textit{AMC 10 10A 2005}) An equiangular octagon has four sides of length $1$ and four sides of length $\sqrt{2}/2$, arranged so that no two consecutive sides have the same length.  What is the area of the octagon?


 $\textbf{(A) } \frac{7}{2}\qquad \textbf{(B) } \frac{7\sqrt{2}}{2}\qquad \textbf{(C) } \frac{5+4\sqrt{2}}{2}\qquad \textbf{(D) } \frac{4+5\sqrt{2}}{2}\qquad \textbf{(E) } 7$


	\item (\textit{AMC 10 10A 2005}) For how many positive integers $n$ does $1+2+\dotsb+n$ evenly divide $6n$? 


 $\textbf{(A) } 3\qquad \textbf{(B) } 5\qquad \textbf{(C) } 7\qquad \textbf{(D) } 9\qquad \textbf{(E) } 11$


	\item (\textit{AMC 10 10A 2005}) Let $S$ be the set of the $2005$ smallest positive multiples of $4$, and let $T$ be the set of the $2005$ smallest positive multiples of $6$. How many elements are common to $S$ and $T$?


 $\textbf{(A) } 166\qquad \textbf{(B) } 333\qquad \textbf{(C) } 500\qquad \textbf{(D) } 668\qquad \textbf{(E) } 1001$


	\item (\textit{AMC 10 10A 2005}) Let $\overline{AB}$ be a diameter of a circle and $C$ be a point on $\overline{AB}$ with $2 \cdot AC = BC$. Let $D$ and $E$ be points on the circle such that $\overline{DC} \perp \overline{AB}$ and $\overline{DE}$ is a second diameter. What is the ratio of the area of $\triangle DCE$ to the area of $\triangle ABD$?


 
                    \begin{figure}[H]
                        \centering
                        \includegraphics[width=0.5\textwidth]{Problem/AMC_10_10A_2005_23_Problem_Figure_1.png}
                    \end{figure}
                    
$\textbf{(A) } \frac{1}{6} \qquad \textbf{(B) } \frac{1}{4} \qquad \textbf{(C) } \frac{1}{3} \qquad \textbf{(D) } \frac{1}{2} \qquad \textbf{(E) } \frac{2}{3}$


	\item (\textit{AMC 10 10A 2005}) For each positive integer $m > 1$, let $P(m)$ denote the greatest prime factor of $m$. For how many positive integers $n$ is it true that both $P(n) = \sqrt{n}$ and $P(n+48) = \sqrt{n+48}$?


 $\textbf{(A) } 0\qquad \textbf{(B) } 1\qquad \textbf{(C) } 3\qquad \textbf{(D) } 4\qquad \textbf{(E) } 5$


	\item (\textit{AMC 10 10A 2005}) In $\triangle ABC$ we have $AB = 25$, $BC = 39$, and $AC = 42$. Points $D$ and $E$ are on $\overline{AB}$ and $\overline{AC}$ respectively, with $AD = 19$ and $AE = 14$. What is the ratio of the area of triangle $ADE$ to the area of the quadrilateral $BCED$?


 $\textbf{(A) } \frac{266}{1521}\qquad \textbf{(B) } \frac{19}{75}\qquad \textbf{(C) } \frac{1}{3}\qquad \textbf{(D) } \frac{19}{56}\qquad \textbf{(E) } 1$


\end{enumerate}